\documentclass[11pt,a4paper]{amsart}
\usepackage[top=2cm,left=2.2cm,right=2.2cm,bottom=2cm]{geometry}
\usepackage[T1]{fontenc}
\usepackage[utf8]{inputenc}
\usepackage{mathptmx}
\usepackage{microtype}
\usepackage[colorlinks,urlcolor=blue!50!black,linkcolor=black]{hyperref}
\usepackage{xcolor}

\pagestyle{plain}
\setlength{\parindent}{0pt}
\setlength{\parskip}{0.45em}
\newcommand{\proj}[1]{\par\smallskip\noindent\textbf{#1}\;}

\begin{document}

\begin{center}
{\large\bfseries Statement of Purpose}\\[2pt]
SLMath Research Membership --- \emph{Algebraic Combinatorics}, Fall 2027\\[2pt]
Maciej Dołęga, Institute of Mathematics, Polish Academy of Sciences (IM PAN), Warsaw
\end{center}

\smallskip

My research interests lie at the intersection of integrable probability, enumerative geometry, and mathematical physics, and the driving force that bridges these areas is algebraic and enumerative combinatorics. More specifically, I focus on the combinatorics of Jack and Macdonald symmetric functions, and on how their hidden combinatorial structure explains phenomena in the the aforementioned fields. Each of my current projects needs exactly the mix of communities that the Fall~2027 semester brings to Berkeley: algebraic combinatorialists working on symmetric functions and integrability, and, through the concurrent program \emph{New Trends in Tropical Geometry}, the tropical and real enumerative geometers. I plan to use my stay to advance the following four projects.

\proj{A. $b$-Hurwitz numbers: towards a geometric theory.}
In~[J8] we introduced \emph{$b$-Hurwitz numbers}: a deformation of Hurwitz numbers by a parameter $b$ measuring non-orientability. Their generating function is a specific partition function of Jack symmetric functions --- one-parameter deformation of the classical Schur symmetric functions, which are recovered at the special case of $b=0$. For generic parameter $b$, they count possibly non-orientable branched coverings of the sphere, weighted by $b^{\text{non-orientable}}$, where the exponent is strictly positive if and only if the source is non-orientable. In particular, this model is a one-parameter interpolation between the classical complex geometry of Hurwitz numbers ($b=0$) and their real counterpart case ($b=1$), which places for the first time Jack symmetric functions within real algebraic geometry. With collaborators I showed that they are computed by refined topological recursion~[J1], and arise from Whittaker vectors for $\mathcal{W}$-algebras~[E3], and in some specific cases they correspond to a partition function of various matrix models, sych as the BGW integrals, and they satisfy the integrable BKP hierarchy~[J5]. This places an integrable model inside real enumerative geometry, bringing both new results and new questions there. Tropical geometers have taken the theory up: Hahn--Markwig found tropical and polynomial structures in the case $b=1$ (\emph{twisted Hurwitz numbers}), and Fesler--Hahn--Karev--Markwig used tropicalization to answer our question on the polynomiality of simple $b$-Hurwitz numbers.

The central open problem I want to attack is geometric: \emph{what is the geometric interpretation of the full multiparametric Jack generating series?} This question is wide open: even for the case of product of length three (known to the algebraic geometers as the Grothendieck dessins d'enfant in the classical $b=0$ case), it corresponds to a longstanding conjecture of Goulden and Jackson from 1996 known as \emph{the $b$-conjecture}, which is still unresolved. My earlier work settled special cases of the $b$-conjecture~[J17,~J18,~J20] and it motivated our work on the closely related two-point multiparametric Jack expansions~[J8], which led to the discovery of the aforementioned $b$-Hurwitz numbers and brought Jack symmetric functions to the field of real enumerative geometry. Our work in progress suggests that a key to success is to understand how combinatorial aspects of certain integrable operators can be related with the tropical aspects of real enumerative geometry. The joint workshops of the two programs, and daily contact with the tropical and real enumerative geometers of the concurrent program (in particular H.~Markwig and her collaborators), make the semester the ideal setting for this question.

\proj{B. Positivity for Jack polynomials.}
Besides the \emph{$b$-conjecture} and its geometric extension described in the previous section, there is another longstanding conjecture in the area of combinatorics of symmetric functions that I find particularly beautiful. Lassalle, motivated by the asymptotic representation theory of the symmetric goup and its connection to free proability theory, conjectured in 2009 that there exist certain universal polynomials, known as \emph{Kerov polynomials for Jack characters} that express the normalized coefficients of Jack polynomials in the power-sum basis in terms of free cumulants, and that these polynomials have non-negative integer coefficients. With H.~Ben~Dali, I found a positive formula for Jack polynomials and Jack characters and proved two conjectures of Lassalle, in particular the integrality of Kerov polynomials for Jack characters~[E4], but their positivity remains open.  

In our work H.~Ben~Dali, we have merged two different combinatorial approaches to the problem: the first one uses the combinatorial maps, which are natural combinatorial objects that appear in the $b$-Hurwitz theory, and the second one uses various operators acting on symmetric functions, and brings a combinatorial understanding of these operators through lattice paths. All the studied operators come from a representation of the Schiffmann--Vasserot algebra $\mathrm{SH}^c$ on symmetric functions, and the new tool that we would like to develop within this project is a general combinatorial interpretation of this representation in the spirit of~[E4,~E1], which should give a structural route to the Kerov positivity problem. Symmetric functions, integrability and combinatorial representation theory are core themes of the program, and its concentration of experts on Macdonald and Jack combinatorics is what this project needs to succeed.

\proj{C. Discrete $\beta$-ensembles at high temperature.}
Discrete $N$-particle systems with Jack-type interaction are discrete counterparts of $\beta$-ensembles and a central object of integrable probability. The asymptotic behaviour of these systems when their size $N$ tends to infinity was recently studied for a fixed $\beta$, and in analogy with the continuous models of $\beta$-ensambles, it does not depend on the deformation parameter $\beta$ when properly scaled. However, the situation in the so-called high-temeperature $\beta\to0$, $\beta N\to\mathrm{const}$, is very different and this asymptoptic behaviour was not well-understood until very recently. With C.~Cuenca we characterized all such systems that have a limit shape in the high-temperature regime in~[J2]. This work is an example where a probabilitic problem is solved by almost only algebraic combinatorial tools: our proof uses the action of Cherednik--Dunkl operators on symmetric functions to compute the moments of the empirical measure, and our combinatorial analysis of this action expresses these models as weighted sum of lattice paths. There are two immediate questions stemming from this work:
\begin{enumerate}
\item[(i)] what is the ``shape'' of this limiting object?
\item[(ii)] what are the fluctuations of the system around the limiting object?
\end{enumerate}
We discovered in a subsequent work~[E2] an interesting \emph{crystallization} phenomenon that does not occur in the world of continuous $\beta$-ensambles in the high-temperature regime: in three specific cases the limit density is a uniform measure supported on intervals with gaps of length one, so that it is uniquely described by the right endpoints of these intervals, that we call \emph{crystals}. Moreover, we have proved that certain classical hypergeometric functions are real-rooted, and the crystal locations coincide with the real zeros of these functions in these three cases. We believe that this is a general phenomenon, and we plan to use the program to prove it in full generality. In particular, we plan to locate the crystals for the whole class as the real zeros of a ``convolution'' of real-rooted entire functions. In the other direction, we plan to use the combinatorial framework developed in~[J2] to study the fluctuations of these systems. The main challenge is to control not only the leading term of the mixed moments, but also the subleading terms. We believe that all these terms are fully controlled by a combinatorial framework of lattice paths, where several interesting statistics are responsible for the whole asymptotic expansion. This project lies at the meeting point of two of the program's focus areas, \emph{integrable systems and dynamical combinatorics} and \emph{geometry of polynomials}. Discussions with S.~Corteel, Ph.~Di~Francesco and L.~Williams (integrable combinatorics, tilings and arctic phenomena) and with C.~Vinzant (real-rootedness) would be particularly valuable.

\proj{D. General-$\beta$ matrix integrals, with J.~Novak.}
With J.~Novak (UCSD) I study the large-$N$ asymptotics of the general-$\beta$ analogues of the Harish-Chandra--Itzykson--Zuber and Brézin--Gross--Witten integrals. These are multivariate Bessel functions; the second also appears in gauge theory as a Gaiotto state. Novak's approach to the case $\beta=2$ uses monotone Hurwitz numbers, and their combinatorial aspects. The missing general-$\beta$ ingredients come from our recent theory: $b$-monotone Hurwitz numbers with their cut-and-join equation and Virasoro constraints~[J5,~J8], and the asymptotics of Jack--Plancherel diagrams~[J14,~J3]. We have already proved a version of the conjectured asymptotics and are now writing the paper. The more ambitious general statement remains challenging and wide open. As far as I am concerned Jonathan Novak plans to take part in the program, so overlapping at SLMath may be decisive for completing the project, as we otherwise have little opportunity to work together in person.

\proj{Contribution to the program.}
I would be happy to give a talk or short tutorial/seminar session on any of the presented topics. I would also be glad to run a working group on Jack and $b$-deformations across combinatorics, integrability and real and tropical enumerative geometry. I have organized similar working groups at IM~PAN and at the Jagiellonian University. I have mentored eight postdocs and co-supervised four PhD students, and would gladly mentor program postdocs and students. Finally, I am the main organizer of the Simons Semester \emph{Combinatorial structures, integrable probability, and topological recursion} at IM~PAN in spring~2027. I will bring the questions and collaborations emerging from it to Berkeley, and help strengthen the links between SLMath and the Polish combinatorics community. I plan to be in residence for the first half of the program, which covers the workshops \emph{Pathways Workshop: Algebraic Combinatorics \& New Trends in Tropical Geometry}, \emph{Introductory Workshop: Algebraic Combinatorics \& New Trends in Tropical Geometry}, and \emph{Recent Developments in Algebraic Combinatorics}. I would also like to mention that my wife, Jacinta Torres, is a mathematician working in combinatorial representation theory and she is also applying for this program, as we are planning to move to Berkeley together. We would be happy to contribute to the program as a couple, and we are both very excited about the opportunity to be part of it.

\end{document}
